\scanpage{98}

$X\in\mathfrak g$ tels que

\[
X\deleted{\unsure{\mathfrak g|j}}V^\beta\subset\sum_{\beta'\geq\beta}V^{\beta'}\qquad\text{pour tt }\beta\in\mathbb Z.
\]

Il suffit que l'on ait

\[
X\in\sum_{\alpha\geq0}\mathfrak g^\alpha,
\]

et c'est aussi nécessaire si $V$ est un module fidèle, car si $X=\sum_{\alpha\in\mathbb Z}X^\alpha$, et si \deleted{$\unsure{\lambda|\alpha}$} $\mu$ est le plus petit des $\alpha$ tels que $X^\alpha\ne0$, i.e. $X=X^\mu+\sum_{\lambda>\mu}X^\lambda,\ X^\mu\ne0$,

\deleted{\unsure{si on avait}} $\exists\,\alpha$ tel. que $X^\mu\cdot V^\alpha\ne0$, \deleted{donc} et comme

\[
X^\mu V^\alpha\subset V^{\alpha+\mu}
\]

il faut qu'on ait $\unsure{v|\nu|\alpha}+\mu\leq\alpha$ i.e. $\mu\leq0$. On voit de $\overline{\mathrm m}$ que si $V$ est fidèle, on a $\mathfrak u=\mathrm{Lie}\,U=\sum_{\alpha>0}\mathfrak g^\alpha$. Par passage à la limite, on en conclut aisément, dans tous les cas

\[
\mathfrak p=\sum_{\alpha\geq0}\mathfrak g^\alpha,\qquad \mathfrak u=\sum_{\alpha>0}\mathfrak g^\alpha.
\]

\deleted{Supposons \unsure{$G$} encore que $G$ soit de t. fini, soit \unsure{un module} fidèle $V$. \unsure{Alors} On a alors}

\deleted{$V$ tel que $P=P_V,\ U=U_V$. On a donc, si $\unsure{P'|Q'}$ est le $\unsure{P'\subset P}$ \unsure{ss-gpe} parabolique de $P'=\underline{\mathrm{Aut}}(V)$, qui stabilise la filtration de $V$, et $U'$ le ss-gpe \unsure{qui induit} qui induit id sur $\mathrm{Gr}(V)$}

\[
\deleted{P=G\cap P',\qquad U=G\cap U'}
\]

\deleted{et par suite on a une imm. fermée}

\[
\deleted{G/P\unsure{\subset|\longleftarrow}G'/P'}
\]
% archive sheet number: 97
